\documentclass[11pt,a4paper]{article}

\usepackage[T1]{fontenc}
\usepackage[utf8]{inputenc}
\usepackage{lmodern}
\usepackage{amsmath,amssymb}
\usepackage{geometry}
\usepackage{hyperref}

\geometry{margin=1in}

\title{Memory-Channel and Mode-Locking Constraints in TRM/TQM V3.0}
\author{
Fabrice Wieser\\
Independent Researcher\\
\url{https://github.com/MagusDraconis/TRM_Cosmology}
}
\date{June 2026\\
V3.0 Review Baseline\\
Release tag: \texttt{trm-v3.0-review-release}}

\begin{document}
\maketitle

\begin{abstract}
This paper isolates two tightly coupled TRM/TQM V3.0 tracks: the memory-channel derivation path and the rational mode-locking closure path \cite{memory_derivation,m3_closure_path}. The goal is to document what is currently tested-effective and structurally constrained, while preserving explicit claim boundaries \cite{trm_status_v3}. The memory term \(\phi^2\lvert\dot{\mu}\rvert\) is currently strongly supported as a minimal effective invariant under guard tests, and the \(m=3\) closure route is currently strongly constrained through RBF hardening \cite{trm_memory_first_principles,trm_rational_band_first_principles}. Neither path is claimed as theorem-level microscopic closure.
\end{abstract}

\section{Introduction}

TRM/TQM V3.0 uses a test-gated strategy: derive candidate structures, stress them under ablations, and retain only forms that survive structural and numerical constraints \cite{trm_status_v3}. Within this strategy, the memory-channel and mode-locking blocks are central because they constrain transport response and closure order, respectively.

This document focuses on:
\begin{enumerate}
  \item memory-channel admissibility and current microscopic support level;
  \item mode-locking closure structure and the current \(m=3\) theorem-path hardening;
  \item explicit failure criteria and remaining first-principles gaps.
\end{enumerate}

\section{Memory-Channel Track}

\subsection{Core object}

The effective memory-channel invariant is currently represented by:
\[
\phi^2\lvert\dot{\mu}\rvert.
\]

\subsection{Current test support}

The memory path is guarded by MEM/TRM/MC blocks, including MC09-MC12 hardening, and currently supports \cite{memory_derivation,trm_memory_first_principles,geodesic_derivation,finsler_action}:
\[
A_{\mathrm{dyn}}\propto\phi
\;\rightarrow\;
A_{\mathrm{dyn}}^2\lvert\dot{\mu}\rvert
\;\rightarrow\;
\phi^2\lvert\dot{\mu}\rvert.
\]

At V3.0, this is interpreted as strong lattice-proxy/microscopic support at effective level, not theorem-level derivation.

\subsection{Claim boundary}

\begin{itemize}
  \item Supported: tested-effective/hypothesis-supported effective invariant path.
  \item Not claimed: theorem-level microscopic uniqueness proof.
\end{itemize}

\section{Rational Mode-Locking and \(m=3\) Closure Track}

\subsection{Core object}

The mode-locking path is organized around closure-family constraints and the inverse rational band \cite{collective_mode_locking,trm_rational_band_first_principles}:
\[
\Omega=\frac{q+3}{q},\qquad \gamma=\frac{1}{\Omega}
\]
\[
\Omega \approx 1.16..1.19,\qquad \gamma \approx 0.84..0.86.
\]

\subsection{Current hardening status}

The Gap-2 path is currently hardened through RBF blocks, with RBF16-RBF23 providing \cite{m3_closure_path,trm_status_v3}:
\begin{enumerate}
  \item connected threshold-region stability,
  \item explicit failure-by-family exclusion,
  \item bounded perturbation stability,
  \item microscopic phase-lattice-energy-based action/tick discriminator reinforcement.
\end{enumerate}

This keeps \(m=3\) as a strongly constrained closure-order candidate under the tested rule family.

\subsection{Claim boundary}

\begin{itemize}
  \item Supported: strongly constrained theorem path for \(m=3\).
  \item Not claimed: theorem-level microscopic closure theorem.
\end{itemize}

\section{Coupled Interpretation: Memory + Mode-Locking}

The two tracks are methodologically coupled:
\begin{enumerate}
  \item the memory-channel block constrains admissible transport-memory structure;
  \item the mode-locking block constrains admissible closure order and discriminator structure.
\end{enumerate}

Together they reduce arbitrary model freedom, but they do not yet establish full microscopic first-principles completion.

\section{Falsification and Failure Criteria}

The current interpretation should be narrowed or rejected if:
\begin{enumerate}
  \item a competing memory invariant satisfies the same admissibility constraints with strictly better bridge behavior and no additional penalties;
  \item \(m=3\) loses constrained minimality under robust, non-ad-hoc perturbation families;
  \item microscopic derivation assumptions force a different leading invariant or closure order under the same structural constraints.
\end{enumerate}

\section{Remaining Gaps}

Open items remain explicit:
\begin{itemize}
  \item theorem-level microscopic closure for \(\phi^2\lvert\dot{\mu}\rvert\),
  \item theorem-level microscopic closure for \(m=3\),
  \item full first-principles derivation bridge between these blocks and broader unified action structure.
\end{itemize}

\section{Conclusion}

In V3.0, the memory-channel and mode-locking tracks are substantially hardened and claim-safe: both are strongly constrained and tested-effective in their current role, but neither is promoted to theorem-level microscopic proof. This status supports critical external review and targeted next-step falsification work, not over-claiming completion.

\bibliographystyle{plain}
\bibliography{references}

\end{document}
